\documentclass[11pt,a4paper]{article}

% --- 核心宏包与中文支持 ---
\usepackage[utf8]{inputenc}
\usepackage{xeCJK}        % 必须使用 XeLaTeX 编译器以支持中文
\usepackage{amsmath, amssymb, amsfonts}
\usepackage{listings}     % 用于代码块显示
\usepackage{xcolor}
\usepackage{geometry}
\usepackage{tcolorbox}    % 教学框样式
\usepackage{booktabs}
\usepackage{tabularx}
\usepackage{setspace}
\usepackage{enumitem}
\usepackage[hidelinks]{hyperref}
\usepackage{bookmark}
\usepackage{fancyhdr}

% --- PDF 书签与元数据 ---
\hypersetup{
    colorlinks=false,
    pdfborder={0 0 0},
    bookmarks=true,
    bookmarksnumbered=true,
    bookmarksopen=true,
    bookmarksopenlevel=2,
    pdfcreator={XeLaTeX},
    pdftitle={02 - Control Structures},
    pdfauthor={黄子扬},
}

% --- 页面美化设置 ---
\geometry{top=0.7in,bottom=0.8in,left=0.6in,right=0.6in}
\onehalfspacing
\tcbuselibrary{skins,breakable}
\setlength{\parindent}{0pt}
\setlength{\parskip}{0.6em}

% --- 代码高亮深度配置 ---
\definecolor{mainblue}{RGB}{0,51,102}
\definecolor{examred}{RGB}{180,0,0}
\definecolor{codebg}{RGB}{248,248,248}
\definecolor{commentgreen}{RGB}{0,128,0}

\lstset{
    backgroundcolor=\color{codebg},
    basicstyle=\ttfamily\footnotesize,
    keywordstyle=\color{blue}\bfseries,
    commentstyle=\color{commentgreen},
    stringstyle=\color{orange},
    breaklines=true,
    showstringspaces=false,
    frame=leftline,
    xleftmargin=2em,
    numbers=left,
    numbersep=5pt,
    numberstyle=\tiny\color{gray},
    language=Python,
    morekeywords={None, True, False, def, return, if, elif, else, for, in, while, break, continue, pass, range, and, or, not}
}

% --- 自定义教学方框 ---
\newtcolorbox{concept}[1]{colback=blue!5, colframe=mainblue, fonttitle=\bfseries, title={{【Concept / 核心概念】#1}}, arc=3pt, breakable}
\newtcolorbox{warning}[1]{colback=red!5, colframe=examred, fonttitle=\bfseries, title={{【Warning / 避坑指南】#1}}, arc=3pt, breakable}
\newtcolorbox{examplebox}[1]{colback=green!2, colframe=green!40!black, fonttitle=\bfseries, title={{【Example / 实战案例】#1}}, arc=2pt, breakable}

% --- 页眉页脚 ---
\pagestyle{fancy}
\fancyhf{}
\fancyhead[L]{\small\bfseries 02 - Control Structures / 控制结构}
\fancyhead[R]{\small\thepage}
\renewcommand{\headrulewidth}{0.5pt}

% --- 文档信息 ---
\title{\Huge \bfseries 02 - Control Structures}
\author{\Large \textbf{哈尔滨工业大学 \quad 黄子扬}}
\date{2026年6月8日}

\begin{document}

\maketitle
\tableofcontents
\newpage

% ======================================================================
\section{Boolean Recap / 布尔值回顾}
% ======================================================================

\begin{concept}{Review from Lecture 01}
    Booleans are one of the data types. They can only be one of two values: \texttt{True} or \texttt{False}. \\
    布尔型只有两个值：\texttt{True}（真）和 \texttt{False}（假）。

    Booleans can result from:
    \begin{itemize}[itemsep=0.3em]
        \item \textbf{Comparison operators}: \texttt{==, !=, <, >, <=, >=} (e.g., \texttt{41 > 42} $\rightarrow$ \texttt{False}).
        \item \textbf{Logical operators}: \texttt{and}, \texttt{or}, \texttt{not} — these act on Booleans and produce Booleans.
    \end{itemize}
    布尔值可以通过比较运算符和逻辑运算符产生。这是条件判断的基础。
\end{concept}

% ======================================================================
\section{Conditionals: If-Statements / 条件判断}
% ======================================================================

\subsection{Basic Logic / 基础逻辑}

\begin{concept}{What is an If-Statement? / 什么是 If 语句？}
    If-statements allow you to execute parts of your code based on whether a condition is met. \\
    If 语句允许程序根据条件是否满足来决定执行哪段代码。

    \textbf{The Rule / 核心规则}：As long as the statement after the \texttt{if} keyword evaluates to \texttt{True}, the \textbf{indented} code block will run. If the statement is \texttt{False}, all the indented parts will be skipped. \\
    只要 \texttt{if} 后面的条件为 \texttt{True}，缩进的代码块就会运行。如果为 \texttt{False}，缩进部分会被完全跳过。
\end{concept}

\subsection{The Full Structure: if / elif / else}

\begin{concept}{Three Keywords / 三个关键字}
    \begin{itemize}[itemsep=0.5em]
        \item \textbf{\texttt{if}}：The first condition to check. / 第一个必检条件。
        \item \textbf{\texttt{elif}} (else if)：Checks additional conditions \textbf{only if} the previous \texttt{if} (or \texttt{elif}) was \texttt{False}. / 仅在前面的条件为 \texttt{False} 时才检查额外条件。
        \item \textbf{\texttt{else}}：Runs if \textbf{none} of the above conditions were \texttt{True}. It is the fallback. / 以上所有条件都不满足时的兜底执行。
    \end{itemize}
\end{concept}

\begin{examplebox}{Price Classification / 价格分级案例}
    课件中通过以下案例展示了完整的 \texttt{if-elif-else} 结构：
\begin{lstlisting}
price = 42
if price > 100:
    print("expensive")          # Runs if price > 100
elif price > 50:                # Checked ONLY if price <= 100
    print("moderate")           # Runs if 50 < price <= 100
else:                           # Runs if all above are False
    print("cheap")              # Runs if price <= 50
# Output: cheap
\end{lstlisting}
\end{examplebox}

\begin{warning}{Execution Order / 执行顺序（重要）}
    In an \texttt{if-elif-else} chain:
    \begin{enumerate}[itemsep=0.3em]
        \item Python checks each condition \textbf{from top to bottom}.
        \item The \textbf{first} condition that is \texttt{True} gets executed.
        \item All subsequent conditions are \textbf{skipped} — even if they would also be \texttt{True}!
    \end{enumerate}
    从上到下检查，第一个为真的条件被执行，其余全部跳过！
\end{warning}

\subsection{Inline Conditionals / 行内条件判断}

\begin{examplebox}{Simple If-Statement / 简单 If 判断}
    课件展示了只需判断单个条件的场景：
\begin{lstlisting}
name = "John"
if len(name) < 3:
    print("your name is short!")    # Won't run because len("John") = 4
\end{lstlisting}
    如果条件为 \texttt{False}，缩进代码块什么都不执行。
\end{examplebox}

\begin{concept}{Ternary Operator (One-line if/else) / 三元运算符}
    Python also supports a one-line conditional expression for simple true/false assignments:
\begin{lstlisting}
# Syntax: value_if_true if condition else value_if_false
result = "short" if len(name) < 3 else "not short"
\end{lstlisting}
    三元运算符适合简单的二选一赋值，不适合复杂逻辑。
\end{concept}

% ======================================================================
\section{Functions / 函数基础}
% ======================================================================

\subsection{Why Functions? / 为什么需要函数？}

\begin{concept}{Building Blocks of Python / Python 的积木块}
    Functions are the building blocks of almost every Python program. They're where the real action takes place. \\
    函数是几乎所有 Python 程序的基本构建块。

    Functions break code into smaller chunks, and are great for:
    \begin{itemize}[itemsep=0.3em]
        \item Tasks that a program uses \textbf{repeatedly} — instead of writing the same code each time, just call the function.
        \item Breaking complex problems into manageable pieces.
        \item Making code more readable and maintainable.
    \end{itemize}
\end{concept}

\begin{concept}{Built-in vs. User-Defined / 内置函数 vs 自定义函数}
    \begin{itemize}[itemsep=0.3em]
        \item \textbf{Built-in functions}：Come built into Python itself (e.g., \texttt{print()}, \texttt{len()}, \texttt{type()}).
        \item \textbf{User-defined functions}：Functions you create yourself to perform specific tasks.
    \end{itemize}
    你已经用过 \texttt{print()} 和 \texttt{len()} 等内置函数，现在你将学会定义自己的函数。
\end{concept}

\subsection{How Function Calls Work / 函数调用的三步流程}

\begin{concept}{The Three Steps of a Function Call / 函数调用的三个步骤}
    \begin{enumerate}[itemsep=0.5em]
        \item \textbf{Call / 调用}：The function is called, and any arguments are passed to the function as input.
        \item \textbf{Execute / 执行}：The function executes, and some action is performed with the arguments.
        \item \textbf{Return / 返回}：The function returns, and the original function call is \textbf{replaced with the return value}.
    \end{enumerate}
\end{concept}

\begin{examplebox}{Understanding len() / 理解 len() 的执行过程}
\begin{lstlisting}
num_letters = len("emlyon")
print(num_letters)          # Output: 6
\end{lstlisting}
    What happens behind the scenes:
    \begin{enumerate}
        \item \texttt{len()} is called with the argument \texttt{"emlyon"}.
        \item The length of the string is calculated — the number 6.
        \item \texttt{len()} returns 6 and replaces the function call with the value.
        \item Python then assigns \texttt{6} to \texttt{num\_letters}. Equivalent to: \texttt{num\_letters = 6}
    \end{enumerate}
\end{examplebox}

\begin{warning}{Calling Errors / 调用错误}
    Functions like \texttt{len()} require exactly one argument. Calling \texttt{len()} without input will raise a \texttt{TypeError}:\\
    \texttt{TypeError: len() takes exactly one argument (0 given)}

    \begin{itemize}[itemsep=0.3em]
        \item Some functions can be called with \textbf{no} arguments.
        \item Some functions can take \textbf{as many arguments as you like}.
        \item \texttt{len()} requires \textbf{exactly one} argument.
    \end{itemize}
\end{warning}

\subsection{Writing Your Own Functions / 定义自己的函数}

\begin{concept}{Two Parts of a Function / 函数的两大部分}
    Every function has two parts:
    \begin{enumerate}
        \item \textbf{Function Signature / 函数签名}：Defines the name of the function and any inputs (parameters) it expects. Starts with the \texttt{def} keyword.
        \item \textbf{Function Body / 函数体}：Contains the code that runs every time the function is used. Must be \textbf{indented}!
    \end{enumerate}
\end{concept}

\begin{examplebox}{Creating Your First Function / 创建你的第一个函数}
\begin{lstlisting}
def multiply(x, y):         # Function signature: name + parameters
    product = x * y         # Function body (indented!)
    return product          # Return statement: sends result back

# Calling the function
result = multiply(2, 4)     # result = 8
print(result)               # 8
\end{lstlisting}
\end{examplebox}

\begin{concept}{Key Points About User-Defined Functions / 自定义函数要点}
    \begin{itemize}[itemsep=0.3em]
        \item \texttt{def} keyword creates a new function. The function is assigned to a variable with the same name as the function name.
        \item Function names become variables, so they must follow the \textbf{same naming rules} as variables.
        \item You call a user-defined function just like any other function: \texttt{function\_name(arguments)}.
        \item \texttt{return} sends a value back. If there is no \texttt{return} statement, the function returns \texttt{None}.
    \end{itemize}
\end{concept}

\subsection{Arguments and Return Values / 参数与返回值}

\begin{concept}{Key Terminology / 关键术语}
    \begin{itemize}[itemsep=0.5em]
        \item \textbf{Parameter / 形参}：A variable in the function signature that defines what input the function expects (e.g., \texttt{x} and \texttt{y} in \texttt{def multiply(x, y)}).
        \item \textbf{Argument / 实参}：The actual value passed to the function when calling it (e.g., \texttt{2} and \texttt{4} in \texttt{multiply(2, 4)}).
        \item \textbf{Return Value / 返回值}：When a function is done executing, it returns a value as output. The function call is replaced by this value.
    \end{itemize}
\end{concept}

% ======================================================================
\section{Scope / 作用域的秘密}
% ======================================================================

\begin{concept}{Local vs. Global / 局部与全局}
    Names (variable names) are unique. Each function has its own \textbf{local scope} with its own set of names. Code outside functions is in the \textbf{global scope}. \\
    变量名是唯一的。函数内部拥有独立的"局部作用域"，函数外部的代码属于"全局作用域"。
\end{concept}

\begin{examplebox}{Scope Demonstration / 作用域演示}
\begin{lstlisting}
x = 2                       # Global x / 全局变量

def func():
    x = 3                   # NEW local x — does NOT change global x!
    print("Inside:", x)     # Prints 3 (local x)

func()                      # Output: Inside: 3
print("Outside:", x)        # Output: 2 (global x unchanged!)
\end{lstlisting}
    \textbf{关键理解}：函数内部的 \texttt{x = 3} 创建了一个\textbf{全新的局部变量}，不会影响全局的 \texttt{x}。
\end{examplebox}

\begin{concept}{Scope Hierarchy / 作用域层级}
    Scopes have a hierarchy. When Python encounters a variable name:
    \begin{enumerate}[itemsep=0.3em]
        \item First, it looks in the \textbf{local scope} (inside the current function).
        \item If not found, it moves up to the \textbf{global scope} (outside all functions).
        \item If still not found, it raises a \texttt{NameError}.
    \end{enumerate}
    Python 查找变量时先从局部找，找不到才去全局找。
\end{concept}

\begin{examplebox}{Scope Lookup Example / 作用域查找示例}
\begin{lstlisting}
x = 5                       # Global

def func():
    y = x + 5               # x not found locally → uses global x (5)
    print("Inside, y =", y) # Output: Inside, y = 10

func()
print("Outside, x =", x)    # Output: Outside, x = 5
# print(y)                  # NameError! y only exists in func()'s local scope
\end{lstlisting}
    \textbf{注意}：函数内部可以\textbf{读取}全局变量，但不能直接修改它（除非用 \texttt{global} 关键字）。
\end{examplebox}

% ======================================================================
\section{Loops / 循环：程序如何做重复工作}
% ======================================================================

\begin{concept}{Why Loops? / 为什么需要循环？}
    Programmers don't want to work hard — they want to work \textbf{smart}. Loops are the tool for that. \\
    程序员不想做重复劳动——循环就是用来解决这个问题的。

    Imagine your boss wants you to print numbers from 1 to 10. Without loops, you'd write \texttt{print()} ten times. With loops, it's just a few lines. \\
    想象老板让你打印 1 到 10。没有循环你得写 10 次 \texttt{print()}，有了循环只需要几行代码。

    Loops help with tasks that are similar and repetitive, from simple (printing numbers) to complex (nested loops with conditionals).
\end{concept}

\subsection{For Loops / For 循环}

\begin{concept}{For Loop Structure / For 循环结构}
    \begin{lstlisting}
for [every_element] in [my_sequence]:
    do something
    \end{lstlisting}
    \begin{itemize}[itemsep=0.3em]
        \item \texttt{every\_element}：A variable created by the loop that takes on each value in the sequence, one at a time.
        \item \texttt{my\_sequence}：The sequence to iterate over (list, tuple, string, range, etc.).
        \item The loop runs once for each item in the sequence.
    \end{itemize}
\end{concept}

\subsubsection{The range() Function / range() 函数}

\begin{concept}{range(start, stop) / range() 详解}
    \texttt{range(1, 11)} produces numbers from \texttt{1} up to (but \textbf{NOT} including) \texttt{11} — i.e., 1, 2, 3, …, 10. \\
    \texttt{range(1, 11)} 产生 1 到 10 的序列，\textbf{不包含}终点 11。

    \textbf{Syntax / 语法}：
    \begin{itemize}[itemsep=0.3em]
        \item \texttt{range(stop)} — starts at 0, goes up to (but not including) stop.
        \item \texttt{range(start, stop)} — starts at start, goes up to (but not including) stop.
        \item \texttt{range(start, stop, step)} — starts at start, increments by step each time.
    \end{itemize}
\end{concept}

\begin{examplebox}{Printing Numbers with for / 用 for 循环打印数字}
\begin{lstlisting}
for i in range(1, 11):
    print(i)                # Prints 1, 2, 3, ..., 10 (each on new line)
\end{lstlisting}
    其中 \texttt{i} 是迭代变量，依次取值为 1, 2, 3, …, 10。
\end{examplebox}

\subsection{While Loops / While 循环}

\begin{concept}{When to Use While? / 什么时候用 While？}
    Used when we only know \textbf{when to stop}, but not \textbf{how many times} to run. \\
    用于只知道"什么时候停"，但不知道要运行几次的情况。

    \textbf{Structure / 结构}：
    \begin{lstlisting}
while [condition]:
    do whatever
    \end{lstlisting}
    Before starting each iteration, the condition is checked. If \texttt{True}, the loop runs. If \texttt{False}, the loop exits.
\end{concept}

\begin{examplebox}{While Loop vs For Loop / While 循环与 For 循环对比}
\begin{lstlisting}
# For loop version — finite, known iterations
for i in range(1, 11):
    print(i)

# While loop version — same result, different approach
counter = 1
while counter < 11:
    print(counter)
    counter = counter + 1       # MUST update counter!
\end{lstlisting}
\end{examplebox}

\begin{warning}{Infinite Loops / 无限循环风险}
    A \texttt{while} loop could run \textbf{forever} if the condition never becomes \texttt{False}! \\
    如果条件永远不为假，\texttt{while} 循环会无休止运行！

    \textbf{Key Differences / For 与 While 的主要区别}：
    \begin{itemize}[itemsep=0.3em]
        \item \textbf{For loop}：Creates a variable to iterate over a given sequence. It is \textbf{finite} — you know exactly how many times it will run.
        \item \textbf{While loop}：Does NOT create an iteration variable for you. You must create and update your own counter. It \textbf{could run forever} if you're not careful.
    \end{itemize}
\end{warning}

\subsection{Control Keywords / 循环控制关键字}

\begin{concept}{break, continue, pass / 三大控制关键字}
    这三个关键字可以改变循环的正常执行流程。
\end{concept}

\begin{examplebox}{Control Keywords Demo / 控制关键字对比}
\begin{lstlisting}
# Example: break — 立即终止整个循环
for i in range(1, 11):
    if i == 5:
        break           # Exit loop completely when i = 5
    print(i)            # Prints: 1, 2, 3, 4 (stops at 5)
\end{lstlisting}

\begin{lstlisting}
# Example: continue — 跳过本次循环剩余部分
for i in range(1, 11):
    if i == 5:
        continue        # Skip the rest of THIS iteration
    print(i)            # Prints: 1, 2, 3, 4, 6, 7, 8, 9, 10 (no 5!)
\end{lstlisting}

\begin{lstlisting}
# Example: pass — 占位符，什么都不做
for i in range(1, 11):
    if i == 5:
        pass            # Do nothing — just a placeholder
    print(i)            # Prints: 1, 2, 3, 4, 5, 6, 7, 8, 9, 10
\end{lstlisting}
\end{examplebox}

\begin{table}[h]
\centering
\renewcommand{\arraystretch}{1.5}
\begin{tabular}{@{}lll@{}}
\toprule
\textbf{Keyword} & \textbf{Effect / 作用} & \textbf{Use Case / 使用场景} \\ \midrule
\texttt{break}    & Aborts the loop entirely / 直接终止整个循环 & 搜索到目标后退出 \\
\texttt{continue} & Skips the rest of the current iteration / 跳过本次循环剩余部分 & 跳过不需要处理的特殊情况 \\
\texttt{pass}     & A placeholder that does nothing / 占位符 & 预留代码位置，防止空代码块报错 \\ \bottomrule
\end{tabular}
\end{table}

% ======================================================================
\section{Summary / 速查表}
% ======================================================================

\begin{table}[h]
\centering
\small
\renewcommand{\arraystretch}{1.3}
\begin{tabular}{@{}ll@{}}
\toprule
\textbf{Structure / 结构} & \textbf{Syntax / 语法} \\ \midrule
If 判断 & \texttt{if condition:} \\
Elif 判断 & \texttt{elif condition:} \\
Else 兜底 & \texttt{else:} \\
定义函数 & \texttt{def name(params):} \\
返回结果 & \texttt{return value} \\
For 循环 & \texttt{for var in sequence:} \\
While 循环 & \texttt{while condition:} \\
终止循环 & \texttt{break} \\
跳过本次 & \texttt{continue} \\
占位符 & \texttt{pass} \\
Range 序列 & \texttt{range(start, stop, step)} \\ \bottomrule
\end{tabular}
\end{table}

\end{document}
